\section{Introduction}
\label{sec:intro}
The useful unit of GPU ownership is rarely an entire device for the entire
lifetime of an application. A model server needs a prompt response at some
moments, a large KV allocation at others, and sustained memory throughput
during decode. An adapter update or an evaluation job may tolerate interruption
while consuming the arithmetic capacity that the server leaves available.
Assigning the same indivisible device to either application hides these
differences. Sharing the device without controlling interference can be worse:
a throughput improvement is not useful if an interactive request misses its
response-time objective.

Nourai and Bastani's \emph{The GPU Needs an Operating System} makes this case
through large-model measurements and a hardware--software research agenda
\cite{Nourai2026GPUOS}. We adopt its central organizing idea: GPU resources
should be granted according to the urgent workload's remaining slack, rather
than treating low average utilization as unconditional permission to run more
work. We develop that idea into an architectural specification for \system{},
covering process state, resource contracts, compiler checkpoints, admission,
revocation, memory leases, and the hardware interface. This is a system-design
paper supported by component experiments. The architecture includes services
that the current prototype does not yet implement; the implementation and
evaluation identify that boundary explicitly.

The resulting abstraction is a \emph{resource lease attached to restartable
work}. A lease specifies which resources a task may consume, how its progress
is committed, and how much revocation delay admission must reserve. A task can
run in a fused device executor, as a sequence of compiler-generated kernel
slices, or in a separate protected CUDA context. These are execution backends
with different costs and guarantees, not interchangeable meanings of
preemption. The architecture chooses among them using a workload-specific
cost model and preserves one accounting interface across the alternatives.

The prototype results expose limits as well as opportunities. Qualified
fused search improves on native cached execution, but matched CLC+PDL
remains faster. Cooperative sharing benefits a streaming proxy, while
whole-step suspension yields little work beside full-model replay.
The OS must therefore be able to reject an elaborate scheduling path
when ordinary co-running is better.

Our contribution is an explicit architectural specification, an execution
contract separating checkpoint safety from resource return, and component
evidence for bounded implementations. The resource-vector and slack-driven
sharing ideas originate in the cited work; new hardware services and a
completed general-purpose implementation are not claimed.

\section{Why GPU Sharing Needs an OS Contract}
\label{sec:motivation}
\subsection{Decode and the roofline}
Consider a matrix with $W$ weights, represented by $s$ bytes per weight, and
a decode batch of $B$ tokens. A weight-only approximation performs $2BW$
floating-point operations while reading $sW$ bytes. Its arithmetic intensity
is therefore
\begin{equation}
 I_{\mathrm{weights}} = \frac{2B}{s}\quad\text{FLOP/byte}.
\end{equation}
For two-byte weights this is $B$; for one-byte weights it is $2B$. KV reads,
activation traffic, routing, and synchronization add costs that the expression
does not capture. The device ridge $R=F_{\mathrm{peak}}/\beta_{\mathrm{mem}}$
compares peak arithmetic throughput and memory bandwidth at a consistent
precision. A memory-bound operation has an idealized arithmetic-utilization
bound $\min(1,I/R)$. A low value is a structural property of the operation,
not evidence that a dispatch queue was implemented poorly.

For a mixture-of-experts layer, the relevant batch is the batch reaching one
expert. Under balanced routing, $E$ experts and $k$ active experts per token
give an average $B_e=Bk/E$. A large global batch can consequently remain a
small local matrix problem. Routing imbalance makes the average a poor
description of the tail. This matters for an OS: scheduling against the
global token count can overestimate the arithmetic reuse of each ready task,
while the latest-finishing expert determines when dependent work is released.

\paragraph{Evidence motivating the design.}
OpenInfer reports roughly 7\% tensor activity for Qwen3.5-397B at a
97.2\% deadline-attainment operating point on four RTX PRO 6000 cards,
and 2.9\% in its decode-only pool~\cite{Nourai2026GPUOS}. Its Kimi K2
experiments report 6--11\% activity at carried loads. These are that paper's
serving measurements, not our RTX 5090 observations. Its streaming control
reaches approximately 99\% of measured bandwidth on 24 of 188 SMs, whereas
its foreground reaches only about half of peak bandwidth. This contrast
motivates jointly choosing a foreground kernel and its allocation.

\subsection{Why a scheduler alone is insufficient}
Batching can increase reuse, but waiting for a larger batch consumes response
time and larger steps can increase token latency. A scheduler cannot remove
that tradeoff by rearranging identical non-preemptible kernels. It needs
additional work that is both useful and compatible with the urgent phase.
Compatibility depends on more than arithmetic intensity: a cache-resident
matrix loop and an operand-streaming GEMM may issue similar tensor work while
causing very different foreground interference.

Nor can the scheduler infer available SMs from low bandwidth utilization.
A poorly sized foreground kernel may need many SMs to generate enough
outstanding memory requests. Taking those SMs away reduces the foreground's
own throughput before the background transfers a single byte. Conversely,
a well-pipelined streaming kernel may saturate bandwidth with fewer workers.
The resource cost of a partition therefore has two parts: the foreground's
loss from a smaller execution allocation and additional interference from
the admitted background workload. The architecture measures these parts
separately whenever a suitable matched control exists.

\subsection{Time, space, and authority}
CUDA stream priorities influence scheduling but do not promise to interrupt
already-running work~\cite{CUDAStreams}. Green contexts constrain spatial
execution resources; MPS coordinates multiple clients; MIG provides a
different isolation boundary on supported devices~\cite{CUDAGreen,NVIDIAMPS}.
None should be described as a software-selected, instruction-level deadline
primitive merely because it changes latency in one workload. Hardware
context switching, compiler checkpoints, and static partitions address
different parts of the problem.

We distinguish three authorities. An application can state urgency and
declare its work. A runtime can admit or withhold future slices and manage
buffers it owns. A driver or device can enforce context and physical-resource
boundaries. The proposed OS API exposes the strongest available guarantee
without silently upgrading it. In particular, a software issue-rate limit is
not a memory-controller reservation, and a CTA launch limit is not a promise
that a particular set of SMs belongs exclusively to that client.

\section{System Architecture and Objects}
\label{sec:architecture}
\subsection{A six-layer decomposition}
We retain the source paper's six-layer decomposition as a design vocabulary
\cite{Nourai2026GPUOS}. L0 provides device signals, context scheduling,
queues, and enforceable resource controls. L1 defines the obligations of a
shareable kernel. L2 dispatches eligible work to execution resources. L3
governs resource consumption and memory capacity. L4 estimates interference
and revocation costs. L5 selects admissions within foreground slack.
Table~\ref{tab:layers} maps this specification to the current artifact.

\begin{table*}[t]
\centering\small
\begin{tabular}{p{.10\textwidth}p{.38\textwidth}p{.43\textwidth}}
\toprule
Layer & Architectural responsibility & Artifact boundary \\
\midrule
L0 & Context switch, bounded pause, resource enforcement & Existing driver switching is measured; no new device pause or memory-QoS interface. \\
L1 & Restartable work and declared resource footprint & Handwritten WMMA checkpoints and strict FX/Triton slices for supported tensor graphs. \\
L2 & Ready-work selection and dispatch & Fused per-frontier device work pool, stream-ordered worker control, and a separate host slice executor. \\
L3 & Capacity leases, issue pacing, cache/copy governance & Bounded arenas and CTA issue pacing; no physical cache or DRAM reservation. \\
L4 & Workload-specific interference and revoke-cost estimates & Measured configuration lookup; no predictor for unobserved workloads. \\
L5 & Deadline-aware admission and revocation & Held-out slack selection and one-foreground IPC admission; no integrated production service controller. \\
\bottomrule
\end{tabular}
\caption{Architectural specification and implemented evidence. The table is
a capability map, not a claim that all six layers form a finished deployed OS.}
\label{tab:layers}
\end{table*}

\subsection{Processes, tasks, and leases}
A \emph{GPU process} is the unit of ownership and accounting. Its identity
names an address-space authority, loaded executable versions, memory
allocations, and outstanding tasks. Trusted tasks may share one CUDA context
and even one fused executor; mutually untrusted tasks require a protection
boundary supplied below the runtime. The API makes this trust mode explicit
rather than treating a common function-pointer table as process isolation.

A \emph{task} names executable work and its dependency edges. Its descriptor
contains a code version, tensor views, input/output access modes, a progress
format, a resource declaration, and a deadline class. A \emph{slice} is a
bounded portion of that task with one well-defined commit boundary. A
\emph{lease} grants resources to one or more eligible slices for a bounded
interval. A lease also specifies what must happen on revocation: retain a
resident allocation, drain and return CTA resources, save to device memory,
or invoke a driver-managed context transition.

\subsection{Resource vectors and accounting}
The proposed lease resource vector is
\begin{equation}
 \mathbf r=(r_{\mathrm{SM}},r_{\mathrm{TC}},r_{\mathrm{L2}},r_{\mathrm{DRAM}},
 r_{\mathrm{cap}},r_{\mathrm{link}},r_{\mathrm{copy}},r_{\mathrm{power}}).
\end{equation}
The coordinates have different units. SM allocation concerns resident
blocks, warp issue, registers, and shared memory. Tensor work concerns
arithmetic issue. Cache and DRAM budgets concern traffic and queueing.
Capacity is bytes held over time plus the cost of reclaiming them. Link and
copy resources include direction and endpoint. Power affects clock rates
shared by otherwise disjoint work. A scalar utilization percentage cannot
substitute for this vector.

Declarations and measurements serve different purposes. A compiler can
report the logical inputs and scratch required by a slice, along with the
compiled register and shared-memory allocation. It cannot derive physical
DRAM traffic exactly from those fields because cache hits, reuse, and other
clients alter that traffic. The accounting service keeps both values, with
explicit labels. If a counter is unavailable, it remains unknown; it is not
filled with an arithmetic estimate and presented as a measurement.

Resource prices are foreground costs, not universal exchange rates. For
foreground class $f$, background declaration $\mathbf r$, and execution
state $x$, write $C_f(\mathbf r,x)$ for the additional completion time. The
state includes phase, batch geometry, context length, device allocation,
and recent interference. Even if individual resource effects are measurable,
their combination need not be additive. Cache eviction can increase DRAM
traffic, and a power change can alter the service time of every slice.
The implemented lookup does not assume a linear price model; a more general
predictor must be validated on resource combinations before admission uses it.

\subsection{Events and executable versions}
The architecture uses release/acquire publication between producers and
consumers in the appropriate memory domain. A producer publishes the data
and progress state before publishing completion. A consumer observes the
completion event before reading those fields. A host-mapped doorbell requires
system-scope ordering; a flag used only between kernels on the device can use
device scope. Polling a volatile field without a publication protocol is not
a replacement for these ordering rules.

Executable versions have lifetimes too. Hot-loaded code, device function
pointers, and plan descriptors remain valid until every possible caller has
retired. Replacing a jump-table entry does not prove that the old module is
unreachable by a CTA that already loaded its pointer. The specification
therefore associates a code epoch with each admitted task and reclaims an
old version only after an epoch-completion condition. The existing fused
executor uses precompiled handlers; general safe hot replacement remains
an architectural service rather than a property established by these tests.

\section{Compiler Contract and Execution Backends}
\label{sec:compiler}
\subsection{Safe points are semantic boundaries}
A compiler cannot make an arbitrary running GPU kernel safely restartable
by inserting a flag read at an arbitrary instruction. Threads may disagree
about which barrier phase they are in; shared memory may hold an unpublished
partial result; or an external consumer may have observed an output that
the resumed task would produce again. A checkpoint is legal only when its
progress and externally visible effects define an unambiguous restart state.
This requirement is stronger than limiting the number of loop iterations.

For a pure pointwise tile, a commit can consist of final stores to a disjoint
output region. For a matrix product, a commit may consist of a partial
accumulator and the completed K extent. A reduction similarly needs its
partial value and reduction interval. The next invocation consumes exactly
that state. The task's output becomes visible to dependent graph nodes only
after all required tiles have committed. Intermediate numerical state may
remain private even when some task-local work is complete.

The proposed general contract admits opaque kernels only if their provider
supplies equivalent safe-point semantics or an acceptable non-preemptible
bound. Unsupported side effects, device allocation, inter-block protocols,
or external collectives cannot be assumed safe merely because the front end
can capture their call site. Rejecting a graph is a valid compiler result.
A silent eager fallback would violate the scheduler's assumed pause contract
and is therefore inappropriate for this backend.

\subsection{Implemented FX-to-Triton lowering}
The current front end captures shape-specialized, pure inference graphs with
FX and also exposes a \texttt{torch.compile} backend~\cite{Reed2022FX}.
It lowers supported operations to Triton kernels and generated pointwise
expressions~\cite{Tillet2019Triton}. It accepts floating CUDA tensors with
checked shape, stride, dtype, and device specializations. Supported graphs
include broadcast arithmetic, common activations, two-dimensional matrix
products and linear layers, contiguous last-dimension sums and means,
bounded-row softmax, and selected no-copy views. Autograd, mutation,
unrecognized operations, and unsupported semantic options are rejected.

The compiler generates a sequence of stages rather than forcing every graph
into a persistent megakernel. Each stage specifies its number of independent
program instances, its reduction chunks, output storage, and any private
partial storage. A matrix-product slice covers a bounded set of output tiles
and a bounded K interval, accumulating in FP32. A sum or mean saves FP32
partials between reduction chunks. An elementwise stage maps a bounded
range of logical elements through broadcasting and physical strides. A
softmax row currently executes as one program and is limited to 8192 elements.

The execution cursor identifies the stage, program offset, and reduction
chunk. It lives in the host executor; partial arithmetic state lives in
device allocations. The executor advances through all program ranges of a
chunk before advancing the chunk, allowing successive grants to change
their program counts without skipping or duplicating output tiles. The
same stream orders those grants. Resuming on another stream requires an
explicit handoff; the prototype rejects an unordered stream change.

Every slice is a finite kernel launch. Once that launch completes, its CTA
resources are returned to the CUDA scheduler. This is the important semantic
difference from a resident worker that sleeps. The cost is additional launch
and synchronization work, more partial-state traffic, and possibly worse
cache reuse. The compiler exports resource and progress metadata so that the
controller can account for these costs. A maximum program count or K extent
is a bound on work geometry, not a guaranteed number of microseconds.

\subsection{Fused dynamic execution}
The earlier fused executor, called FuseOS in the retained measurements,
handles a different execution opportunity: changing work within a bounded
model invocation. In output-driven search, scores determine the next live
branch set, selected histories, and ready tiles. The model architecture alone
does not fix those choices. Stable request identifiers and bounded storage
plans separate changing logical ownership from physical reservations.

The evaluated Qwen3-0.6B path stores an immutable ancestor KV region and two
private-history banks. Each frontier copies selected histories into the
alternate bank, appends KV, and swaps banks. Precompiled handlers combine
normalization and projection, gate/up and activation, and residual work.
A cooperative persistent invocation executes all 28 layers and the vocabulary
head for one frontier. Device work counters assign ready tiles within a
phase; barriers order dependent phases. Host branch scoring remains part of
the complete-request measurement. This is not a process-lifetime serving
kernel with a GPU-owned external request queue.

The same arithmetic handlers support ordinary grids, programmatic dependent
launch, and cluster launch control. Shape-specific graphs cache encountered
widths and depths. A fixed-topology control accepts an updated device
configuration without recapture. These controls matter architecturally:
the OS must preserve good fused arithmetic and plan reuse even when it
selects an execution backend other than persistence. It must not create an
artificially weak graph baseline to justify a resident runtime.

\paragraph{A process-lifetime resident service.}
The architecture also permits a resident executor to accept requests from
a host-visible ring, perform sampling and stopping decisions on the device,
and publish tokens through an output ring. This is the process-lifetime
execution model demonstrated for Llama-3.1-8B by OpenInfer, which also
assigns selected CTAs background roles at phase boundaries
\cite{Nourai2026GPUOS}. Such a service needs request-slot generations,
backpressure, cancellation, and a rule for retiring every output consumer.
A full ring must stop admissions without blocking token delivery for accepted
requests. A cancelled request retires only after its current phase commits
and references to its KV are released. Our frontier-level executor provides
bounded model execution but does not implement this complete service loop.

\subsection{Three meanings of suspension}
\system{} assigns an explicit suspension mode to a backend. A resident
cooperative worker saves progress and stops useful issue while retaining its
CTA allocation. A finite sliced backend commits and returns from the kernel,
releasing CTA allocation before resumption. A separate-context backend relies
on driver-managed switching, whose saved physical execution state is opaque
to the application. All can preserve logical task progress, but only the
second exposes an application-controlled kernel boundary with resource return.